% GENERATED FILE. Edit canonical lessons, then run npm run generate:books.
% canonical-source-hash: fnv1a64:baf023db4109dd6d
% canonical-lessons: FR-C391-epargne, FR-C391-epargner, FR-C391-investir, FR-C391-investissement, FR-C391-taux

\chapter{Mettre de côté}
\label{ch:fr-mettre-de-cote}
\hlchaptermodality{391}

\begin{chapteropening}
\textbf{By the end of this chapter:} \emph{I can talk about saving and investing: savings, saving money, investing, an investment and the interest rate.}

Complete the last lesson of chapter 391: I can talk about saving and investing: savings, saving money, investing, an investment and the interest rate.
\end{chapteropening}

\section[l'épargne]{l'épargne — the savings, saving}

\label{lesson:FR-C391-epargne}



\begin{quote}
Before the new one: say the French for the expense, then the French for the contribution.
\end{quote}

\subsection*{You'll want to know: l'épargne}
\textbf{l'épargne} — \textquotedblleft{}the savings, saving\textquotedblright{}.

\begin{quote}
Ils ont mis leur épargne sur un compte à la banque.
\end{quote}

\begin{quote}
\emph{They put their savings in a bank account.}
\end{quote}

Feminine. Formal; in speech, \emph{les économies}.

\begin{cousinweb}
\emph{épargne} is formed from \emph{épargner}, inherited through Old French \emph{espargnier} from a Frankish \emph{sparanjan} \textquotedblleft{}to spare\textquotedblright{} (EtymDB).
\end{cousinweb}

\begin{grammarlens}[title={What you've built}]
One more word for telling, explaining and describing at B1.
\end{grammarlens}

\subsection*{Your turn}
\begin{itemize}
\raggedright
  \item \emph{Say it:} \emph{l'épargne}
  \item \emph{Say it:} \emph{l'épargne}, once more
  \item \emph{Say it:} a sentence of your own with \emph{l'épargne}
  \item \emph{Recall:} say the French for the expense, then the French for the contribution, then say \emph{l'épargne} again
\end{itemize}

\begin{tcolorbox}[breakable,colback=teal!4,colframe=teal!35!black,title={Before you move on}]
What does l'épargne mean? (\textquotedblleft{}The savings, saving\textquotedblright{}.) Say it once more.
\end{tcolorbox}
\section[épargner]{épargner — to save up (money); to spare}

\label{lesson:FR-C391-epargner}



\begin{quote}
Before the new one: say the French for the contribution, then the French for the savings.
\end{quote}

\subsection*{You'll want to know: épargner}
\textbf{épargner} — \textquotedblleft{}to save up (money); to spare\textquotedblright{}.

\begin{quote}
Elle épargne un peu chaque mois.
\end{quote}

\begin{quote}
\emph{She saves a little every month.}
\end{quote}

Also to spare someone: \emph{épargner quelqu'un}.

\begin{cousinweb}
\emph{épargner} is inherited from Old French \emph{espargnier} (EtymDB), from a Frankish \emph{sparanjan} \textquotedblleft{}to spare\textquotedblright{}.
\end{cousinweb}

\begin{grammarlens}[title={What you've built}]
One more word for telling, explaining and describing at B1.
\end{grammarlens}

\subsection*{Your turn}
\begin{itemize}
\raggedright
  \item \emph{Say it:} \emph{épargner}
  \item \emph{Say it:} \emph{épargner}, once more
  \item \emph{Say it:} a sentence of your own with \emph{épargner}
  \item \emph{Recall:} say the French for the contribution, then the French for the savings, then say \emph{épargner} again
\end{itemize}

\begin{tcolorbox}[breakable,colback=teal!4,colframe=teal!35!black,title={Before you move on}]
What does épargner mean? (\textquotedblleft{}To save up (money); to spare\textquotedblright{}.) Say it once more.
\end{tcolorbox}
\section[investir]{investir — to invest}

\label{lesson:FR-C391-investir}



\begin{quote}
Before the new one: say the French for the savings, then the French for to save up.
\end{quote}

\subsection*{You'll want to know: investir}
\textbf{investir} — \textquotedblleft{}to invest\textquotedblright{}.

\begin{quote}
Il a investi son argent dans un appartement.
\end{quote}

\begin{quote}
\emph{He invested his money in a flat.}
\end{quote}

A \emph{-ir} verb like \emph{finir}: \emph{nous investissons}.

\begin{cousinweb}
\emph{investir} is a learned borrowing from Latin \emph{investīre} \textquotedblleft{}to clothe, to cover\textquotedblright{} (EtymDB; Lewis \& Short), built on \emph{vestis} \textquotedblleft{}clothing\textquotedblright{}; the money sense came later, probably through Italian.
\end{cousinweb}

\begin{grammarlens}[title={What you've built}]
One more word for telling, explaining and describing at B1.
\end{grammarlens}

\subsection*{Your turn}
\begin{itemize}
\raggedright
  \item \emph{Say it:} \emph{investir}
  \item \emph{Say it:} \emph{investir}, once more
  \item \emph{Say it:} a sentence of your own with \emph{investir}
  \item \emph{Recall:} say the French for the savings, then the French for to save up, then say \emph{investir} again
\end{itemize}

\begin{tcolorbox}[breakable,colback=teal!4,colframe=teal!35!black,title={Before you move on}]
What does investir mean? (\textquotedblleft{}To invest\textquotedblright{}.) Say it once more.
\end{tcolorbox}
\section[l'investissement]{l'investissement — the investment}

\label{lesson:FR-C391-investissement}



\begin{quote}
Before the new one: say the French for to save up, then the French for to invest.
\end{quote}

\subsection*{You'll want to know: l'investissement}
\textbf{l'investissement} — \textquotedblleft{}the investment\textquotedblright{}.

\begin{quote}
Une bonne formation est un bon investissement.
\end{quote}

\begin{quote}
\emph{Good training is a good investment.}
\end{quote}

Masculine.

\begin{cousinweb}
\emph{investissement} is built from \emph{investir} and \emph{-ment}. \emph{investir} comes from Latin \emph{investīre} \textquotedblleft{}to clothe\textquotedblright{}, built on \emph{vestis} \textquotedblleft{}clothing\textquotedblright{} (Lewis \& Short).
\end{cousinweb}

\begin{grammarlens}[title={What you've built}]
One more word for telling, explaining and describing at B1.
\end{grammarlens}

\subsection*{Your turn}
\begin{itemize}
\raggedright
  \item \emph{Say it:} \emph{l'investissement}
  \item \emph{Say it:} \emph{l'investissement}, once more
  \item \emph{Say it:} a sentence of your own with \emph{l'investissement}
  \item \emph{Recall:} say the French for to save up, then the French for to invest, then say \emph{l'investissement} again
\end{itemize}

\begin{tcolorbox}[breakable,colback=teal!4,colframe=teal!35!black,title={Before you move on}]
What does l'investissement mean? (\textquotedblleft{}The investment\textquotedblright{}.) Say it once more.
\end{tcolorbox}
\section[le taux]{le taux — the rate}

\label{lesson:FR-C391-taux}



\begin{quote}
Before the new one: say the French for to invest, then the French for the investment.
\end{quote}

\subsection*{You'll want to know: le taux}
\textbf{le taux} — \textquotedblleft{}the rate\textquotedblright{}.

\begin{quote}
Le taux d'intérêt est de trois pour cent.
\end{quote}

\begin{quote}
\emph{The interest rate is three per cent.}
\end{quote}

The same form in the plural. \emph{Le taux de chômage} — the unemployment rate.

\begin{cousinweb}
\emph{taux} comes from Old French \emph{tauxer} \textquotedblleft{}to fix a price\textquotedblright{} (EtymDB), from Latin \emph{taxāre} \textquotedblleft{}to appraise, to value\textquotedblright{} (Lewis \& Short).
\end{cousinweb}

\begin{grammarlens}[title={What you've built}]
One more word for telling, explaining and describing at B1.
\end{grammarlens}

\subsection*{Your turn}
\begin{itemize}
\raggedright
  \item \emph{Say it:} \emph{le taux}
  \item \emph{Say it:} \emph{le taux}, once more
  \item \emph{Say it:} a sentence of your own with \emph{le taux}
  \item \emph{Recall:} say the French for to invest, then the French for the investment, then say \emph{le taux} again
\end{itemize}

\begin{tcolorbox}[breakable,colback=teal!4,colframe=teal!35!black,title={Before you move on}]
What does le taux mean? (\textquotedblleft{}The rate\textquotedblright{}.) Say it once more.
\end{tcolorbox}
